\section{Main chapters}

\subsection{Chapter 1: Problem Statement and Theoretical Analysis of Market Shocks}

\subsubsection{Purpose of the chapter:}

To describe in detail the problem of modeling market shocks, justify the relevance of the study, formulate research hypotheses and determine the main parameters of the model.

\subsubsection{Content:}

\begin{enumerate}
    \item Theoretical background of market shocks:
    \begin{enumerate}
        \item Definition of a market shock and its main characteristics (sudden change in price, liquidity, trading volume)
        \item Review of the causes of shocks: external economic events, regulatory changes, technological failures
        \item Analysis of the effects of shocks on financial markets based on empirical studies (see Danielsson, 2005; Brunnermeier, 2009)
    \end{enumerate}
    \item The relevance of modeling market shocks:
    \begin{enumerate}
        \item The rationale for the need for modeling to predict and prevent systemic crises.
        \item Advantages of the agent-based approach compared to traditional economic modeling techniques (Epstein & Axtell, 1996; Tesfatsion, 2003).
    \end{enumerate}
    \item Formulation of research objectives:
    \begin{enumerate}
        \item To investigate the dynamics of the model without shocks and establish baselines.
        \item To introduce different scenarios of market shocks followed by comparative analysis.
        \item Conduct parametric analysis of the influence of various factors (number of traders, intensity of events) on the behavior of the system.
    \end{enumerate}
\end{enumerate}

\subsubsection{Conclusion of the chapter:}

Brief summary of the main points, formulated hypotheses and objectives of the research, preparation of ground for description of modeling methodology in the next chapter.

\subsection{Chapter 2: Modeling methodology and description of the simulator}

\subsubsection{Purpose of the chapter:}

To present the technical and methodological background of the study, to describe the architecture of the simulator used, and to justify the choice of parameters and experiment scenarios.

\subsubsection{Content:}

\begin{enumerate}
    \item Description of the simulator:
    \begin{enumerate}
        \item AgentBasedModel/traders.py - all traders and their parameters
        \item AgentBasedModel/simulator.py - Simulator object and its initialization
        \item AgentBasedModel/exchange.py - Exchange objects, assets and their initialization
        \item AgentBasedModel/extra/events.py - what events can be triggered on the stock exchange during trading
        \item AgentBasedModel/visualization/ - built-in visualization tools
    \end{enumerate}
    \item Customization of the experimental environment:
    \begin{enumerate}
        \item Description of the procedure for creating and setting up the environment (installing dependencies, setting up the Jupyter kernel).
        \item Determination of initial values of model parameters and justification of their choice based on literature data (Lux & Marchesi, 1999).
    \end{enumerate}
    \item Experimental scenarios:
    \begin{enumerate}
        \item Formulation of the baseline scenario (model without exposure to shock events).
        \item Design of experimental scenarios modeling market shocks (e.g., sudden decrease in liquidity, sudden change in demand).
        \item Methodology of data collection and processing, and a description of the statistical techniques used (e.g. regression analysis, time series analysis).
    \end{enumerate}
\end{enumerate}

\subsubsection{Conclusion of the chapter:}

Conclusions on the feasibility of the chosen modeling approach and methodology, preparation for the experimental study.

\subsection{Chapter 3: Experimental study and analysis of results}

\subsubsection{Purpose of the chapter:}

To conduct a series of experiments using the developed simulator, analyze the impact of market shocks on the system behavior and compare the results with theoretical expectations.

\subsubsection{Content:}

\begin{enumerate}
    \item Description of the experiments:
    \begin{enumerate}
        \item Presentation of the basic simulation results in the absence of shock events.
        \item Description of the experiments with the introduction of different shock scenarios:
        \begin{enumerate}
            \item Sudden decrease in liquidity.
            \item A sudden increase in trading volume when the price falls.
        \end{enumerate}
        \item Detailed description of changes in key parameters and methodology of experiments.
    \end{enumerate}
    \item Data analysis:
    \begin{enumerate}
        \item Presentation of experimental results in the form of graphs, tables and statistical measures.
        \item Comparison of the system behavior under the baseline scenario and scenarios with shocks.
        \item Discussion of the identified patterns, comparing with the findings of previous studies (see Lux & Marchesi, 1999; Samanidou et al., 2007).
    \end{enumerate}
    \item Discussion of the results:
    \begin{enumerate}
        \item Interpretation of the results obtained in the context of the hypotheses posed.
        \item Analysis of the limitations of the model and possible reasons for discrepancies with theoretical expectations.
        \item Suggestions for further improvement of the model and research methodology.
    \end{enumerate}
\end{enumerate}

\subsubsection{Conclusion of the chapter:}

Summary conclusions on the conducted experiments, evaluation of the effectiveness of agent-based modeling for the analysis of market shocks, formulation of recommendations for further research.